\documentclass[11pt,a4paper]{article}
\usepackage[margin=1in]{geometry}
\usepackage{amsmath,amssymb,booktabs,graphicx,hyperref,siunitx,cleveref}
\usepackage[T1]{fontenc}
\usepackage[dvipsnames]{xcolor}
\hypersetup{colorlinks=true,linkcolor=blue!50!black,citecolor=blue!50!black,urlcolor=blue!50!black}
\usepackage[numbers,sort&compress]{natbib}
\usepackage{authblk}

\author[1]{Triparno Bandyopadhyay}
\author[2]{Rajarajeswari Ramasamy}
\author[2]{Nishitha TP}

\affil[1]{Nayanta University, Bavdhan, Pune 411021, Maharashtra, India.}
\affil[2]{Department of Physics and Nanotechnology, SRM Institute of
Science and Technology, Kattankulathur, Tamil Nadu 603203, India.}

\title{Supplementary notes to \\
\emph{Background Expectation for the
LZ High-Energy Event}}
\date{\today}

\renewcommand{\thesection}{S\arabic{section}}
\renewcommand{\theequation}{S\arabic{equation}}
\renewcommand{\thefigure}{S\arabic{figure}}

\newcommand{\keVee}{\ensuremath{\mathrm{keV_{ee}}}}
\begin{document}
\maketitle

\noindent
These notes collect material that was developed in the course of the
analysis and is not needed to follow its argument. None of it bears on
the central result, which is that the published characterisation of the
electron-recoil response does not determine the tail probability on which
the background expectation depends. We record it because each item either
answers a question a reader may reasonably ask of the main text, or
documents a calculation whose conclusion is quoted there without its
derivation. Section and equation numbers of the main text are referred to
as they appear there.

\section{Contours of constant \(N_q\) in the observable plane}
\label{sup:geom}

Recombination conserves \(N_q\), so it displaces an event along a contour
of constant total quanta. Since \(\mathrm{S1c}=g_1N_{ph}\) and
\(\mathrm{S2c}=g_2N_e\), that contour is
\begin{align}
    \label{eq:locus}
    \frac{\mathrm{S1c}}{g_1}+\frac{\mathrm{S2c}}{g_2}=N_q
    \qquad\Longleftrightarrow\qquad
    \mathrm{S2c}=g_2\!\left(N_q-\frac{\mathrm{S1c}}{g_1}\right),
\end{align}
a straight line of slope \(-g_2/g_1=-312.5\) in the
\(\{\mathrm{S1c},\mathrm{S2c}\}\) plane. In the
\(\{\mathrm{S1c},\log_{10}\mathrm{S2c}\}\) plane in which the bands are
drawn, the same locus is \(\log_{10}[g_2(N_q-\mathrm{S1c}/g_1)]\), a curve
whose local slope
\begin{align}
    \label{eq:locusslope}
    \frac{\mathrm{d}\log_{10}\mathrm{S2c}}{\mathrm{d}\,\mathrm{S1c}}
    =-\frac{g_2/g_1}{\ln 10\;\mathrm{S2c}}
\end{align}
steepens from \(-0.0032\)~dex/phd at the nominal recombination of a
parent of \(69.84\)~\keVee to \(-0.0148\)~dex/phd at the event, and
diverges as \(\mathrm{S1c}\to g_1N_q\). These loci are the
contours of constant energy drawn in grey in Figs.~2 and~4 of the Letter,
where their curvature is directly visible.

The distinction matters when a vertical separation is read off the band.
The bands of the Letter are evaluated taking flat energy spectra along
lines of fixed \(\mathrm{S1c}\), so the reference population at fixed
\(\mathrm{S1c}\) spans many parent energies: at
\(\mathrm{S1c}=540\)~phd the band median corresponds to a parent of
\(89.8\)~\keVee rather than the \(69.84\)~\keVee of the event
inversion. A vertical displacement therefore mixes recombination with
parent energy, whereas a displacement along \cref{eq:locus} isolates
recombination at fixed \(N_q\). The inference of the Letter is a
two-dimensional unbinned profile likelihood ratio and uses neither
projection; we record the distinction because the vertical separation is
quoted elsewhere as a measure of separation from electron-recoil
backgrounds.

\section{The low-energy edge of the acceptance}
\label{sup:floor}

The main text adopts \(0.62\)~\keVee for the lower edge of the
electron-recoil acceptance of the WS RoI and states that three plausible
treatments span \(2.2\%\) in the acceptance width. The derivation is as
follows.

The ceiling is obtained by solving \(Q_y(E)\,E=409.5\) with the tuned
charge-yield model, giving \(20.02\)~\keVee at
\(Q_y=20.46\)~e\(^-\)/keV. The floor cannot be obtained that way, for two
reasons.

First, the yield parameterisation is not usable at sub-keV energies.
Charge yield is bounded above by the case in which no ion recombines,
\begin{align}
    \label{eq:qymax}
    Q_y^{\max}=\frac{1}{W\,(1+N_{ex}/N_i)}=62.6~\mathrm{e^-/keV},
\end{align}
whereas the tuned model tends to \(C\,m_2=85.3\)~e\(^-\)/keV as
\(E\to0\), that is \(m_2=84.9\)~e\(^-\)/keV before the density correction
\(C=1.0051\), exceeding the ceiling by \(36\%\). The two cross at
\(E=0.97~\mathrm{keV_{ee}}\):

\begin{center}
\begin{tabular}{rrrrc}
\toprule
\(E\) [keV] & \(Q_y(E)\) & \(N_e=Q_yE\) & \(N_i=E/[W(1+N_{ex}/N_i)]\) & \(N_e>N_i\)?\\
\midrule
0.193 & 84.6 &  16.3 &   12.1 & yes\\
0.26  & 82.9 &  21.6 &   16.3 & yes\\
0.50  & 78.0 &  39.0 &   31.3 & yes\\
1.0   & 70.6 &  70.6 &   62.6 & marginal\\
5.0   & 44.1 & 220   &  313   & no\\
20.0  & 20.5 & 409   & 1251   & no\\
70.0  & 17.7 & 1237  & 4379   & no\\
\bottomrule
\end{tabular}
\end{center}

\noindent Below about \(1\)~\keVee the parameterisation returns more
electrons than there are ions, and is being used outside the range over
which it was tuned.

Second, the WS RoI is three cuts rather than one. Taking them together,
\(\mathrm{S1c}>3\)~phd requires \(N_{ph}>27.2\) and
\(\mathrm{S2}>645\)~phd requires \(N_e>18.7\) --- the latter binding more
tightly than the S2c floor's \(16.3\) --- so \(N_q>45.9\) and
\(E>0.62~\mathrm{keV_{ee}}\). At that \(N_q\) one has \(N_i=38.9\) and
\(N_{ex}=7.0\), and the two requirements can be met only for
\(r\simeq0.52\), whereas the capped model gives full ionisation there and
\(\mathrm{S1c}=0.78\)~phd. Sub-keV electron recoils therefore enter the
RoI only on a recombination fraction well above the value the capped model
returns, and \(0.62\)~\keVee is a boundary of the acceptance rather
than a threshold that typical events cross.

\begin{center}
\begin{tabular}{lrr}
\toprule
treatment & floor [\keVee] & width [\keVee]\\
\midrule
S2c cut alone, uncapped yield   & 0.193 & 19.80\\
S2c cut alone, full ionisation  & 0.261 & 19.73\\
all three cuts jointly          & 0.621 & 19.37\\
\bottomrule
\end{tabular}
\end{center}

\noindent The spread in \(\mathrm{d}N/\mathrm{d}E\) across the three is
\(2.2\%\). The first is listed only to show what the parameterisation
returns if its sub-keV extrapolation is taken at face value; discarding it
leaves \(1.9\%\).

\section{Bounds on the width-variation sensitivity}
\label{sup:width}

The main text notes that a width known to better than \(+6.2\%\) would
keep the within-family contribution below a tenth of the quoted
background, and that \(+16.6\%\) is where it reaches the whole of it. The corresponding bounds, obtained by folding the varied
response through the parent spectrum, are

\begin{center}
\begin{tabular}{lrr}
\toprule
if \(\sigma_{N_e}\) is high by at most & \(N^{\mathrm{ER}}_{\mathrm{leak}}\) below & fraction of \(0.0106\)\\
\midrule
\(+2\%\)  & \(3.4\times10^{-4}\) & \(3\%\)\\
\(+5\%\)  & \(7.8\times10^{-4}\) & \(7\%\)\\
\(+10\%\) & \(2.7\times10^{-3}\) & \(25\%\)\\
\(+16\%\) & \(9.5\times10^{-3}\) & \(90\%\)\\
\bottomrule
\end{tabular}
\end{center}

\noindent Only upward variations are tabulated; downward variations reduce
the leakage, so these are the conservative side. These bounds constrain
the within-family sensitivity alone. They do not bear on the
identifiability result of the main text, every member of whose families
carries the same \(\sigma_{N_e}\).

\section{The double-vacancy populations}
\label{sup:dec}

The Letter identifies the \(^{124}\)Xe and \(^{125}\)I double-vacancy
decays as the populations whose projection into the tail carries a
systematic it does not propagate. We examined both externally. An external
reconstruction of these populations rests on extrapolations that no
published measurement constrains, and what follows establishes two bounds
and one non-identifiability rather than a determination of the response.
That is why it is recorded here and not in the main text.

These decays produce more light and less charge than equivalent-energy
\(\beta\) decays because the energy is distributed among multiple
low-energy Auger--Meitner electrons, raising the ionisation density and
hence recombination~\cite{Recomb}.

\paragraph{\(^{124}\)Xe.} With \(T_{1/2}=1.09\times10^{22}\)~yr and a KK
fraction of \(64.8\%\) from the collaboration's own
measurement~\cite{Xe124}, the exposure contains \(787\) \(2\nu\)DEC
decays, of which \(510\) are KK captures at \(64.3\)~keV. The rate
calculation can be checked against the same paper's statement that
\(2.3\%\) of the signal lies below the region of interest of that
analysis: that fraction gives \(18\) events against the \(21.0\pm6.3\) of
Table~I of the Letter.

The KK mode is not strongly charge-suppressed. In \(88\%\) of K-captures a
\(28\)~keV X-ray carries energy \(\mathcal{O}(100)~\mu\)m from the primary
site~\cite{Recomb}, so the deposit comprises three independent
recombination sites. Only the both-Auger channel,
\((1-f_K)^2\simeq1.6\%\) or \(\sim\!7\) events, with \(f_K\) the K-shell
fluorescence yield, is a single dense deposit; with the measured box-size
ratio \(a_{\rm L}/a_\beta=0.84\pm0.03\) in the Thomas--Imel
framework~\cite{ThomasImel,Recomb} this gives
\(Q_y/Q_y^\beta\geq0.774\) and a contribution of \(5\times10^{-5}\)
counts. This contribution is negligible.

Taking \(a_{\rm KK}\geq a_{\rm L}\) is the conservative choice. Scaling
the box size with the electron range would give substantially larger
suppression and is not supported by the data: the measured
\(a_{\rm L}/a_\beta=0.84\) corresponds to a range ratio of \(0.10\), so
the box size is far less sensitive to track structure than the range. The
KK double-Auger cascade is moreover dominated by \(\sim\!23\)~keV KLL
electrons, which are more \(\beta\)-like than the sub-keV L cascade.

\paragraph{\(^{125}\)I.} The mechanism is the Letter's own: its Fig.~4
caption records that the spectral features in the ER band above
\(25\)~\keVee come predominantly from \(^{125}\)I and
\(^{133}\)Xe and appear slightly below the \(\beta\)-decay ER band due to
increased recombination from Auger cascades.

\(^{125}\)I decays by electron capture to \(^{125}\)Te, followed by a
\(35.49\)~keV nuclear transition internally converted in \(93\%\) of
decays. That energy lies just above the xenon K edge at
\(34.56\)~keV~\cite{XCOM}, where the mean free path is
\(\simeq0.14\)~mm, so it deposits locally in essentially every decay and
produces a second vacancy; the capture-shell label refers to the first
vacancy only. The K, L and M modes deposit \(67.3\), \(40.4\) and
\(36.5\)~keV, all above the \(20.02\)~\keVee at which a
mean-recombination electron recoil meets the S2c ceiling. The
\(8.9\pm2.7\) \(^{125}\)I events of Table~I therefore cannot be ordinary
electron recoils of the corresponding energies: they are a population
whose \emph{mean} charge yield is suppressed.

\begin{center}
\begin{tabular}{lrrrrr}
\toprule
shell & \(Q_y/Q_y^\beta\) & \(\langle N_e\rangle\) & \(\sigma_{N_e}\) &
\((\langle N_e\rangle-409.5)/\sigma_{N_e}\) & \(P(N_e\le409)\)\\
\midrule
K & 0.773 & 916 & 194 & 2.61 & \(8.3\times10^{-4}\)\\
L & 0.783 & 562 &  99 & 1.54 & \(2.8\times10^{-2}\)\\
M & 0.785 & 517 &  87 & 1.24 & \(7.7\times10^{-2}\)\\
\bottomrule
\end{tabular}
\end{center}

\noindent Here
\(\langle N_e\rangle=(Q_y/Q_y^\beta)\,Q_y(E)\,E\) is evaluated at each
mode's own energy, and \(\sigma_{N_e}\) is the full width including the
binomial term, with \(\omega\) taken from Eq.~4 of the Supplement at that
\(N_q\): it is not a constant, falling from \(0.0456\) at \(N_q=4973\) to
\(0.0369\) at \(N_q=2697\). The implied recombination fractions,
\(1-\langle N_e\rangle/N_i\), are \(0.783\), \(0.778\) and \(0.774\);
these are not the \(Q_y/Q_y^\beta\) ratios, which they resemble
numerically and order in the opposite sense, and each exceeds the nominal
\(\bar{r}\) at that energy by the double-vacancy enhancement. The Gaussian-equivalent
displacements \(\Phi^{-1}(1-P)\) are \(3.1\), \(1.9\) and \(1.4\), larger
than \((\langle N_e\rangle-409.5)/\sigma_{N_e}\) because the distribution is
positively skewed and its low-side tail correspondingly thinner than
Gaussian.

It is important not to overstate this. The L and M modes reach the WS RoI
through \(1.9\) and \(1.4\) Gaussian-equivalent standard deviations,
probabilities of order \(10^{-2}\) to \(10^{-1}\); they are not a probe of
the \(P\simeq10^{-7}\) regime that governs the ordinary electron-recoil
leakage. What places them in the search region is the shift of the mean,
not a rare fluctuation. Only the K mode, at \(3.1\), approaches the
displacement at issue, and only its S1c interval brackets the event at
\(540.1\)~phd: \(\mathrm{S1c}>504\)~phd follows from \(N_e\le409\), and
\(\mathrm{S1c}<549\)~phd from \(N_e\geq0\), that is
\(\mathrm{S1c}\leq N_qg_1\). The L and M modes, at \(N_q=2988\) and
\(2697\), land at \(\mathrm{S1c}\simeq285\)--\(330\) and
\(253\)--\(298\)~phd, inside the WS RoI but well outside the comparison
region. Weighted by capture fraction they supply \(\sim\!92\%\) of the
\(^{125}\)I population in the WS RoI and the K mode the remaining
\(8\%\) --- not because the K mode is absent from its S1c range, but
because reaching the ceiling from a mean of \(916\) demands a much larger
downward fluctuation than the L and M modes require from \(562\) and
\(517\). The implied total is \(\sim\!10^{3}\) decays.

The size of the K-mode population is not recoverable from outside, and the
sensitivity runs parallel to the identifiability result of the main text.
Assigning the L mode the single-vacancy L-capture ratio measured in
Ref.~\cite{Recomb} (\(0.877\)) rather than the double-vacancy treatment
moves the fraction above \(\mathrm{S1c}=504\)~phd from \(8\%\) to
\(82\%\), and the inferred total \(^{125}\)I decay count by an order of
magnitude. Neither assignment is a measurement of the \(^{125}\)I
response. In the absence of a measured double-vacancy response
distribution, modelling choices available externally change the inferred
population in the event's S1c region by a factor of ten --- which is the
regime in which Ref.~\cite{Recomb} states that the Thomas--Imel framework
cannot predict the size of the effect a priori.


\section{An attempted reconstruction of the published band}
\label{sup:band}

Figure~2 of the Letter draws the electron-recoil band as ``the calibrated
NEST model, taking flat energy spectra and evaluating along lines of fixed
S1c''. We implemented that construction, in order to express the deviations
of \cref{sup:geom} against the published band rather than against the
response at a single energy, and did not succeed. We record the attempt
because the main text's quantities are defined at fixed energy partly in
consequence.

At fixed \(\mathrm{S1c}\) the photon number is fixed, and since
recombination conserves quanta, \(N_q=N_{ph}+N_e\). For a flat energy
spectrum the constraint collapses the energy integral exactly,
\begin{align}
    \label{eq:bandcollapse}
    p(N_e\,|\,N_{ph})\;\propto\;\int\!\mathrm{d}E\;f_E(N_e)\,
    \delta\!\left(\tfrac{E}{W}-N_{ph}-N_e\right)
    \;=\;W\,f_{E^\ast}(N_e),\qquad E^\ast=(N_{ph}+N_e)\,W,
\end{align}
so each \(N_e\) is evaluated at the one energy that makes the quanta add up.

The reconstruction reproduces the published \emph{median} to better than
\(0.011\) in \(\log_{10}\mathrm{S2c}\) at every \(\mathrm{S1c}\), and two
independent digitisations of the figure agree with one another to
\(\sim\!0.005\). The \emph{width} does not reproduce:

\begin{center}
\begin{tabular}{rrrrr}
\toprule
\(\mathrm{S1c}\) [phd] & \(E\) at the median [\keVee] &
published \(10\)--\(90\%\) [dex] & reconstructed & ratio\\
\midrule
100 & 17 & 0.134 & 0.148 & 0.91\\
200 & 33 & 0.159 & 0.177 & 0.90\\
300 & 48 & 0.172 & 0.214 & 0.81\\
500 & 82 & 0.230 & 0.291 & 0.79\\
540 & 90 & 0.246 & 0.301 & 0.82\\
\bottomrule
\end{tabular}
\end{center}

\noindent The disagreement is about \(10\%\) at \(17\)~\keVee
and \(20\%\) above \(48\), with the step falling where the electron-recoil
calibration of Ref.~\cite{WS2024} ends. Two candidate causes sit inside our
own assumptions and we cannot separate them from public information. The
range of the ``flat energy spectra'' is not stated; truncating it at
\(\sim\!100\)~\keVee reproduces the published width at
\(\mathrm{S1c}=540\) almost exactly, \(0.250\) against \(0.246\), while
changing nothing below \(\mathrm{S1c}=400\), so a residual \(10\%\) persists
at low \(\mathrm{S1c}\) that the truncation does not touch. Effects we omit
--- Fano smearing of \(N_q\) and the binomial draw of \(N_i\) --- would widen
our band further and so cannot account for the gap, but the unknown
truncation narrows it, and no claim can be made in either direction.

Two things this does \emph{not} indicate. The reconstruction of the response
model itself is verified independently: across the 24 scan energies it
reproduces the \(N_q\), mean \(N_e\) and standard deviation that NEST returns
to \(0.002\%\), \(0.003\%\) and \(0.54\%\), and the variance and skew-normal
placement match the source line for line. And the apparent symmetry of the
published band about its median is an artefact of the logarithmic axis: in
\(N_e\) both bands are positively skewed, with
\((N_{90}-N_{50})/(N_{50}-N_{10})\) rising with \(\mathrm{S1c}\) in each, so
the skewness treatment is not a separate discrepancy from the width.

The scripts, the digitised curves and a fuller account are in the archive, as
\texttt{analysis/band.py} and \texttt{notes/band-reconstruction.md}.

\bibliographystyle{JHEP}
\bibliography{LZ}

\end{document}
